% ============================================================
% shared/theme-werkstatt.tex
% Visuelles Theme für die Werkstatt-Serie (Band 0 und ggf. Folgebände)
% Voraussetzung: xcolor, tcolorbox (most), listings geladen
% ============================================================

% ── Farben ──────────────────────────────────────────────────
\definecolor{WerkstattBlau}{HTML}{1F5F8B}
\definecolor{WerkstattHellblau}{HTML}{EAF4FA}
\definecolor{WerkstattGruen}{HTML}{2E7D32}
\definecolor{WerkstattHellgruen}{HTML}{EEF7EF}
\definecolor{WerkstattOrange}{HTML}{C96A00}
\definecolor{WerkstattHellorange}{HTML}{FFF3E3}
\definecolor{WerkstattGrau}{HTML}{F3F4F6}
\definecolor{WerkstattDunkelgrau}{HTML}{444444}
\definecolor{CodeHintergrund}{HTML}{F7F7F7}

% ── Basis-Einstellungen für alle Boxen ──────────────────────
\tcbset{
  enhanced, breakable,
  boxrule=0.6pt, arc=2mm,
  left=3mm, right=3mm, top=2mm, bottom=2mm,
  before skip=8pt, after skip=8pt
}

% ── Didaktische Boxen ────────────────────────────────────────
\newtcolorbox{aufgabeBox}{
  colback=WerkstattHellorange, colframe=WerkstattOrange,
  fonttitle=\sffamily\bfseries, title=Die Aufgabe}

\newtcolorbox{problemBox}{
  colback=WerkstattGrau, colframe=WerkstattDunkelgrau,
  fonttitle=\sffamily\bfseries, title=Warum es nicht sofort klappt}

\newtcolorbox{wissenBox}{
  colback=WerkstattHellblau, colframe=WerkstattBlau,
  fonttitle=\sffamily\bfseries, title=Was wir dafür lernen müssen}

\newtcolorbox{loesungBox}{
  colback=white, colframe=WerkstattGruen,
  fonttitle=\sffamily\bfseries, title=Schritt-für-Schritt-Lösung}

\newtcolorbox{merksatzBox}{
  colback=WerkstattHellgruen, colframe=WerkstattGruen,
  fonttitle=\sffamily\bfseries, title=Merksatz}

\newtcolorbox{fehlerBox}{
  colback=red!4, colframe=red!55!black,
  fonttitle=\sffamily\bfseries, title=Typische Fehler}

\newtcolorbox{uebungBox}{
  colback=white, colframe=WerkstattBlau,
  fonttitle=\sffamily\bfseries, title=Übungen}

\newtcolorbox{quizBox}{
  colback=WerkstattGrau, colframe=WerkstattBlau,
  fonttitle=\sffamily\bfseries, title=Mini-Quiz}

\newtcolorbox{haubeBox}{
  colback=WerkstattHellblau, colframe=WerkstattBlau,
  fonttitle=\sffamily\bfseries, title=Unter der Haube,
  borderline west={2mm}{0pt}{WerkstattBlau}}

\newtcolorbox{robotikBox}{
  colback=WerkstattHellgruen, colframe=WerkstattGruen,
  fonttitle=\sffamily\bfseries,
  title=Transfer zur Robotik und zu ROS2,
  borderline west={2mm}{0pt}{WerkstattGruen}}

% ── Code-Stile ───────────────────────────────────────────────
\lstdefinestyle{werkstattcode}{
  backgroundcolor=\color{CodeHintergrund},
  basicstyle=\ttfamily\small,
  breaklines=true,
  columns=fullflexible,
  frame=single,
  rulecolor=\color{WerkstattDunkelgrau},
  framerule=0.3pt,
  xleftmargin=2mm, xrightmargin=2mm,
  aboveskip=8pt, belowskip=8pt,
  showstringspaces=false,
  extendedchars=true
}
\lstdefinestyle{bash}{style=werkstattcode, language=bash}
\lstdefinestyle{python}{style=werkstattcode, language=Python}
\lstdefinestyle{text}{style=werkstattcode}
\lstdefinestyle{dockerfile}{style=werkstattcode}
\lstdefinestyle{yaml}{style=werkstattcode}

% ── Kapitel-Makros ───────────────────────────────────────────
\newtcolorbox{analogie}[1][]{
  colback=WerkstattHellblau,
  colframe=WerkstattBlau,
  fonttitle=\sffamily\bfseries,
  title=Analogie,#1}

\definecolor{ProjektFarbe}{HTML}{E8F5E9}
\definecolor{ProjektRahmen}{HTML}{2E7D32}

\newtcolorbox{projektBox}{
  enhanced, breakable,
  colback=ProjektFarbe, colframe=ProjektRahmen,
  boxrule=1pt,
  fonttitle=\sffamily\bfseries,
  title={\texttt{robotik-uno-q} -- Projektstand}}

\newcommand{\kapitelziel}[1]{%
  \begin{tcolorbox}[
    colback=white, colframe=WerkstattBlau,
    fonttitle=\sffamily\bfseries, title=Kapitelziel]
  #1
  \end{tcolorbox}}

\newcommand{\kapitelstruktur}{%
  \section*{Arbeitsweise in diesem Kapitel}
  Jedes Kapitel folgt demselben Muster:
  \begin{enumerate}
    \item Wir starten mit einer konkreten Aufgabe.
    \item Wir prüfen, warum die Aufgabe nicht sofort lösbar ist.
    \item Wir klären, welches Wissen fehlt.
    \item Wir erarbeiten genau dieses Wissen.
    \item Wir lösen die Aufgabe Schritt für Schritt.
    \item Wir übertragen das Gelernte auf Robotik und ROS2.
  \end{enumerate}}
